% =============================================================================
%  JOURNAL PAPER (extension of the UKCI 2026 conference paper)
%
%  Working title:
%    Forecast-Driven Robust Allocation of Critical-Care Surge Capacity
%    in NHS England
%
%  Core contribution: a forecast-to-decision framework in which
%  physics-informed regional demand forecasts drive a REALISTIC, SCALABLE
%  robust optimisation for critical-care surge capacity. The optimisation
%  is the contribution; the forecaster is infrastructure (summarised and
%  cited, not re-derived).
%
%  Positioning:
%    * vs UKCI conference paper : conference = proof of concept (7 regions,
%      snapshot demand, small LP). Journal = operational realism.
%    * vs MSAGAT-Net (Elsevier) : that is forecasting (graph attention);
%      this is decision-making. Forecasting is cited, not the novelty.
%    * vs Shams Eddin & El Hajj (2025): NHS England, critical-care/MV,
%      PINN-estimated flow demand, cost-asymmetric forecasts, equity,
%      exact-vs-metaheuristic at scale.
%
%  TARGET JOURNAL (decided): Health Care Management Science (Springer Nature).
%    - Final submission uses the Springer Nature LaTeX template
%      (`\documentclass[sn-mathphys-num]{sn-jnl}`, sn-jnl.cls). Download from
%      springernature.com/gp/authors/campaigns/latex-author-support .
%    - VERIFY Springer OA coverage for Coventry (oa.lib@coventry.ac.uk):
%      HCMS is hybrid; Coventry's participation in the Springer-Jisc
%      transformative agreement was not confirmed on the library page.
%      Fallback if not covered: EJOR / Computers & OR (Elsevier hybrid,
%      Coventry-covered).
%    - This skeleton keeps `article` so it compiles during drafting; swap to
%      sn-jnl before submission.
% =============================================================================
\documentclass[11pt]{article}

\usepackage[T1]{fontenc}
\usepackage[utf8]{inputenc}
\usepackage{amsmath,amssymb,mathtools}
\usepackage{graphicx}
\usepackage{booktabs}
\usepackage[hidelinks]{hyperref}
\usepackage[capitalise,noabbrev]{cleveref}
\usepackage{geometry}
\geometry{margin=1in}

\title{Forecast-Driven Robust Allocation of Critical-Care Surge Capacity
       in NHS England}
% TODO confirm affiliations. ORCIDs: AmirHosein Sadeghimanesh 0000-0002-6945-3118;
% Matthew England 0000-0001-5729-3420 (Centre Director) -- both Research Centre for
% Computational Science and Mathematical Modelling, Coventry. Fei He, Petra A. Wark:
% confirm Coventry centre/department.
\author{Michael Ajao-Olarinoye \and Abiola Babatunde \and Vasile Palade \and
        AmirHosein Sadeghimanesh \and Fei He \and Petra A. Wark \and
        Matthew England}
\date{Draft -- \today}

\begin{document}
\maketitle

\begin{abstract}
% TODO. ~200 words. Lead with the DECISION problem (where to place surge
% capacity and how to route critical-care patients under demand
% uncertainty), state that physics-informed regional forecasts drive a
% realistic robust optimisation at trust/ICB scale, list the realism
% upgrades (flow demand, correlated scenarios, equity, facility-opening,
% true CVaR/DRO), and the exact-vs-metaheuristic comparison. Close with
% the headline operational result.
\end{abstract}

% =============================================================================
\section{Introduction}\label{sec:intro}
% Problem: NHS critical-care surge planning under demand uncertainty.
% Gap: forecasting papers (incl. our own) optimise accuracy, not decisions;
%   existing allocation papers are not NHS critical-care / not forecast-driven
%   / not at trust scale.
% Contributions (decision-led):
%   (1) forecast-to-decision framework with patient-flow demand;
%   (2) correlated demand scenarios from the joint forecast distribution;
%   (3) trust/ICB-scale robust allocation with facility-opening + equity;
%   (4) exact MILP vs metaheuristic comparison at operational scale;
%   (5) NHS England case study with stress-tested decision value.

% =============================================================================
\section{Related Work}\label{sec:related}
% REUSE + EXPAND docs/paper/sections/02_related_work.tex (currently unused).
% Threads: (a) forecast-to-decision / predict-then-optimise; (b) robust &
% two-stage health-resource location-allocation; (c) ICU/bed allocation &
% metaheuristics; (d) spatiotemporal/physics-informed epidemic forecasting
% (cite MSAGAT-Net, STAN, HOIST, PINNs). Position vs Shams Eddin & El Hajj.

% =============================================================================
\section{Problem Setting and Data}\label{sec:data}
% NEW. Spatial scope: ICB (42 Integrated Care Boards). Demand UNIT =
% ADULT CRITICAL CARE BEDS OCCUPIED (MV beds do not exist publicly below
% NHS region; switch made deliberately and stated).
%   - Demand: per-ICB daily adult critical-care occupancy, NHS England UEC
%     SitRep "Critical Care and General & Acute Beds", trust->ICB, Mar 2020-.
%   - Capacity: per-ICB adult critical-care beds OPEN/available (same SitRep)
%     -> REAL measured capacity, replaces the 1.05*peak proxy. Realism gain.
%   - Geography: ICB centroids (ONS Open Geography Portal); distance/travel
%     matrix (great-circle * detour).
%   - Population: ICB pops via ONS LSOA->ICB lookup + MYE (equity + pop-prop
%     baseline).
%   - Splits: chronological by wave/season; held-out test.

% =============================================================================
\section{Forecasting Module}\label{sec:forecast}
% REUSE / infrastructure (summarise, cite conference paper + ajao2025hybrid +
% MSAGAT-Net as a backend option). Per-ICB forecaster of adult critical-care
% occupancy -> quantile/scenario forecasts. INCLUDE the per-origin bias
% recalibration so the forecaster is accurate across ALL horizons (the
% conference model over-predicts long horizons). Keep short.
% TODO FIGURE (deferred from the conference paper for space): a TikZ
% architecture diagram of the PinnGRU+recal forecaster -- frozen PINN-SEIRD
% encoder (state net U_NN + parameter net X_NN) -> augmented feature window
% -> two-layer GRU -> per-horizon quantile head with the level-and-trend
% anchor -> per-origin recalibration (bias step + persistence blend). The
% conference paper has room only for the high-level pipeline (its Fig. 1);
% the journal version should show the per-layer architecture explicitly.

% =============================================================================
\section{Demand Scenario Generation}\label{sec:scenarios}
% NEW -- the key realism upgrade.
%  5.1 Patient-flow / length-of-stay demand (bed-day, not weekly snapshot).
%  5.2 Correlated scenarios: empirical or Gaussian-copula coupling of the
%      regional forecast residuals (replaces co-monotone q90 stacking).

% =============================================================================
\section{Robust Allocation Model}\label{sec:model}
% NEW / EXTENDED.
%  - first-stage: surge placement + facility-opening (binary x_j at scale);
%  - second-stage recourse: local service, transfers, unmet;
%  - true CVaR or distributionally-robust objective (not the degenerate
%    single-tail penalty of the conference version);
%  - equity term (e.g. max regional shortage ratio, linearised);
%  - budget, capacity, travel-time constraints.

% =============================================================================
\section{Solution Methods}\label{sec:solvers}
% NEW at scale. Exact MILP (CBC/Gurobi) vs metaheuristics (NSGA-II, GA, SA).
% Justified ONLY at trust/ICB scale where exact solution becomes expensive.
% Report solution quality, robustness, and runtime.

% =============================================================================
\section{Case Study and Experimental Design}\label{sec:casestudy}
% NHS England, 42 ICBs. Wave-based split. Operating point + sweeps
% (budget, travel cap, equity weight, robustness level). Stress-test
% instances where demand exceeds capacity.
% BASELINES:
%   Allocation (core): status quo / population-proportional /
%     demand-proportional / greedy shortage-first / deterministic MILP /
%     robust MILP (CVaR-DRO, proposed) / metaheuristics (NSGA-II,GA,SA) /
%     perfect-foresight oracle (upper bound).
%   Forecasting (infrastructure): seasonal-naive, ARIMA, GRU/LSTM, XGBoost,
%     recalibrated PINN-SEIRD; optional MSAGAT-Net backend.
%   Published comparison: Shams Eddin & El Hajj (2025), nearest optimisation
%     result -- re-implement key elements if feasible.

% =============================================================================
\section{Results}\label{sec:results}
% Scale results; robustness; equity trade-offs; exact-vs-metaheuristic
% (quality + runtime); forecast-quality sensitivity; stress-test decision
% value (unmet-demand reduction, not just transfer efficiency).

% =============================================================================
\section{Discussion}\label{sec:discussion}
% Realism achieved vs remaining limitations; policy implications for NHS
% surge planning; generalisability.

% =============================================================================
\section{Conclusion}\label{sec:conclusion}
% Summary + future directions.

% \bibliographystyle{elsarticle-num}
% \bibliography{references}

\end{document}
